% Generated draft: requires mathematical and source review.
\documentclass[11pt]{article}
\usepackage[margin=25mm]{geometry}
\usepackage{fontspec}
\setmainfont{Latin Modern Roman}
\usepackage{amsmath,amssymb,mathtools,unicode-math}
\setmathfont{Latin Modern Math}
\usepackage{enumitem,needspace,xurl,hyperref}
\setlength{\parindent}{0pt}
\setlength{\parskip}{5pt}
\setlength{\emergencystretch}{3em}
\newcommand{\N}{\mathbb{N}}
\newcommand{\Z}{\mathbb{Z}}
\newcommand{\Q}{\mathbb{Q}}
\newcommand{\R}{\mathbb{R}}
\newcommand{\C}{\mathbb{C}}
\newcommand{\F}{\mathbb{F}}
\newcommand{\A}{\mathbb{A}}
\newcommand{\PP}{\mathbb{P}}
\newcommand{\E}{\mathbb{E}}
\newcommand{\eps}{\varepsilon}
\newcommand{\dd}{\,\mathrm{d}}
\begin{document}
\section*{UnsolvedMath ID: 6600003}
\textbf{Source problem code:} catalogue entry
\subsection{Definitions and mathematical statement}
Let $(V,E)$ be the graph associated with a Penrose tiling, with undirected nearest-neighbor edges and counting measure on $V$. Write
\[
\ell^2(V)=\left\{\psi:V\to\mathbb C:\sum_{v\in V}|\psi(v)|^2<\infty\right\}.
\]
Define the graph operator
\[
(H\psi)(v)=\sum_{w:\{v,w\}\in E}(\psi(w)-\psi(v)).
\]
The graph has uniformly bounded degree, so this is a bounded self-adjoint operator on all of $\ell^2(V)$. Its sign convention is adjacency minus degree, the negative of the nonnegative combinatorial Laplacian.

Determine its spectrum
\[
\sigma(H)=\{\lambda\in\mathbb C:H-\lambda I\text{ has no bounded inverse on }\ell^2(V)\}.
\]
The spectrum is real by self-adjointness; no candidate spectral set is supplied. The graph convention needs to be fixed: using tiling vertices joined by tile sides is a usual interpretation, but different Penrose graph models need not define the same operator.

The supplied summary reports infinitely many finitely supported eigenfunctions and corresponding integrated-density-of-states jumps for Penrose graph models, citing arXiv:2209.01443, while leaving a complete spectral description unresolved. Infinitely many such eigenfunctions does not itself assert infinitely many distinct eigenvalues. The summary also mentions four mutually-locally-derivable Penrose graph models and the Ammann--Beenker setting; those are related models, not silently identified with the one operator in this question.
\subsection{Short English statement}
Determine the spectrum of the Penrose-tiling graph operator that sums neighboring value differences. Preserve its negative-Laplacian sign convention and distinguish the intended graph from related tiling models.
\subsection{Sources}
\begin{itemize}
\item \emph{Open Problems and Conjectures Related to the Theory of Mathematical Quasicrystals} (2016), source-order Problem 3 in D. Damanik's ``Quantum Mechanics and Quasicrystals'' contribution, \url{https://arxiv.org/abs/1604.06280}; accessed 2026-07-29, source-TeX extraction.
\item arXiv:2209.01443, supplied reference for locally supported eigenfunctions, multiplicity bounds, and integrated-density-of-states jumps. The record gives no authors or article title for this reference.
\item Local registry \texttt{data/unsolvedmath/problems.json}, permanent UnsolvedMath ID $6600003$, source code catalogue entry, \emph{D. Damanik: Quantum Mechanics and Quasicrystals --- Problem}. Dataset attribution: Ulam AI and the UnsolvedMath Contributors.
\item The supplied rights note reports available arXiv source TeX and a CC BY journal version, with the arXiv item-specific reuse license still NEEDS\_REVIEW.
\end{itemize}
\end{document}
